\chapter{Hordeolum (Stye) and Chalazion}
\index{Hordeolum (Stye) and Chalazion}
\label{sec:Hordeolum (Stye) and Chalazion}

\section{Overview}
A hordeolum (stye) is an acute bacterial infection of the eyelid glands, most commonly caused by \emph{Staphylococcus aureus}. A chalazion is a sterile lipogranulomatous inflammatory lesion that results from obstruction of the meibomian gland duct. Both conditions present with eyelid swelling and tenderness, but they differ in etiology, clinical course, and management.

\section{Classification}
\begin{description}[font=\normalfont\bfseries,leftmargin=3em,labelwidth=2.5em]
  \item[Cause] Infection (hordeolum) or duct obstruction with granulomatous inflammation (chalazion)
  \item[Organ System] Eyes
  \item[Pathophysiology] Hordeolum: bacterial infection of Zeis (external) or meibomian (internal) glands. Chalazion: blocked meibomian gland with retained sebaceous secretions triggering a lipogranulomatous inflammatory response.
  \item[Duration] Acute (hordeolum); subacute to chronic (chalazion)
  \item[Onset] Sudden (hordeolum); gradual (chalazion)
  \item[Mode of Transmission] Autoinoculation from skin flora
  \item[Clinical Course] Hordeolum: painful pustule that may spontaneously drain. Chalazion: painless, firm nodule that persists or slowly enlarges.
  \item[Severity] Mild to moderate
  \item[Extent of Disease] Localized
  \item[Age Group] All ages; more common in adults
  \item[Epidemiology] Very common; hordeolum accounts for a large proportion of ophthalmology outpatient visits; recurrence is frequent.
  \item[ICD-11 Code] 9A00.0 (hordeolum), 9A00.1 (chalazion)
\end{description}

\section{Causes}
Hordeolum is caused by bacterial infection of eyelid glands, most often \emph{Staphylococcus aureus}. Chalazion arises from obstruction of the meibomian gland duct without active infection, leading to retained lipid secretions and a subsequent granulomatous inflammatory reaction.

\section{Risk Factors}
\begin{itemize}[noitemsep]
  \item Blepharitis (chronic eyelid margin inflammation)
  \item Rosacea and seborrheic dermatitis
  \item Poor eyelid hygiene
  \item Contact lens wear
  \item Diabetes mellitus
  \item Previous hordeolum or chalazion
\end{itemize}

\section{Signs and Symptoms}
\begin{itemize}[noitemsep]
  \item Localized eyelid swelling and erythema
  \item Tenderness and pain (more pronounced in hordeolum)
  \item Palpable nodule on the eyelid
  \item Crusting along the eyelid margin
  \item Watering or mild discharge
  \item Blurred vision if large enough to distort the cornea
\end{itemize}

\section{Progression and Stages}
Hordeolum typically progresses over a few days to a pustule that may spontaneously drain and resolve within 1--2 weeks. Chalazion develops slowly over weeks, forming a firm, painless nodule that may persist for months or rarely resolve spontaneously. Recurrent or atypical lesions require biopsy to rule out sebaceous carcinoma.

\section{Complications}
\begin{itemize}[noitemsep]
  \item Recurrence of hordeolum or chalazion
  \item Internal hordeolum extending into orbital cellulitis (rare)
  \item Scarring or eyelid deformity after surgical drainage
  \item Reduced quality of life
  \item Emotional and psychological effects
\end{itemize}

\section{Diagnosis}
\begin{itemize}[noitemsep]
  \item Clinical examination with slit-lamp biomicroscopy
  \item Eversion of the eyelid to inspect the tarsal conjunctiva
  \item Lab tests (not routinely needed)
  \item Imaging tests (not routinely needed)
\end{itemize}

\section{Prevention}
\begin{itemize}[noitemsep]
  \item Maintain good eyelid hygiene with regular cleaning
  \item Avoid sharing cosmetics or eye tools
  \item Treat underlying blepharitis or rosacea
  \item Regular check-ups
\end{itemize}

\section{Treatment Options}
\begin{itemize}[noitemsep]
  \item Warm compresses applied 3--4 times daily for 10--15 minutes
  \item Topical antibiotics (bacitracin or erythromycin) for hordeolum
  \item Incision and curettage for persistent chalazion
  \item Intralesional corticosteroid injection for chalazion
  \item Lifestyle changes
\end{itemize}

\section{Living with It}
\begin{itemize}[noitemsep]
  \item Follow treatment plan
  \item Regular appointments
  \item Adjust daily habits
  \item Support groups
  \item Keep learning
\end{itemize}

\section{Outlook}
Hordeolum usually resolves completely with conservative management or antibiotic therapy within 1--2 weeks. Chalazion may persist longer but responds well to warm compresses or minor surgical intervention. Recurrence occurs in a minority of patients, particularly those with underlying blepharitis or rosacea.
